\chapter{Cauchy's Theorem, Growth and Mellin Inversion}

\begin{center}\itshape
A contour integral vanishes when the contour winds around nothing outside the domain.\\
Growth on the boundary controls growth inside, even on an infinite strip,\\
and a sum of copies at every scale is a Dirichlet series in disguise.
\end{center}

Volume~5 proved Cauchy's theorem on star-shaped sets, and left the general statement: the integral of a holomorphic function over a cycle is zero whenever the cycle winds around no point outside the domain. Volume~7 proves that a simple closed curve has an inside and an outside, and with it the winding number of a continuous curve. This chapter proves Dixon's form of Cauchy's theorem, which uses that winding number and Liouville's theorem. Then come growth and size. The Phragm\'en--Lindel\"of principle extends the maximum principle to unbounded regions and gives the first non-trivial bound for $\zeta(\frac12+it)$; Hadamard's factorisation writes an entire function of finite order from its zeros. The Riemann mapping theorem, which sends every simply connected domain onto a disc, is proved in Volume~9. We end with the Mellin transform, Fourier analysis on the multiplicative group, which turns sums into products of Dirichlet series and is the tool of choice for the asymptotic analysis of algorithms.

\section{Cauchy's theorem in its final form}

A \emph{cycle} is a formal sum $\gamma=\gamma_1+\dots+\gamma_k$ of closed piecewise $C^1$ paths, with $\int_\gamma f=\sum_j\int_{\gamma_j}f$, $\gamma^*=\bigcup_j\gamma_j^*$ and $n(\gamma,a)=\sum_jn(\gamma_j,a)$. The boundary of
an annulus, the outer circle counterclockwise plus the inner circle clockwise, is a cycle but not a single closed path. A cycle $\gamma$ in an open set $\Omega$ is \emph{homologous to zero in $\Omega$} if
$n(\gamma,a)=0$ for every $a\notin\Omega$: it does not wind around any hole of $\Omega$.

\begin{theorem}[Cauchy's theorem, homology form]\label{thm:globalcauchy}
Let $\Omega$ be open, $f$ holomorphic on $\Omega$, and $\gamma$ a cycle in $\Omega$ homologous to zero in $\Omega$. Then
\begin{enumerate}
\item[(a)] $\displaystyle f(z)\,n(\gamma,z)=\frac1{2\pi i}\int_\gamma\frac{f(w)}{w-z}\dd w$ for every $z\in\Omega\setminus\gamma^*$;
\item[(b)] $\displaystyle\int_\gamma f=0$.
\end{enumerate}
\end{theorem}
\begin{proof}
This is John Dixon's proof (1971). Define $g$ on $\Omega\times\Omega$ by $g(z,w)=\frac{f(w)-f(z)}{w-z}$ for $w\ne z$ and $g(z,z)=f'(z)$. It is continuous: away from the diagonal this is clear, and if
$\overline D(z_0,r)\subset\Omega$, then for $z,w\in D(z_0,r)$
\[
g(z,w)=\int_0^1f'\bigl(z+t(w-z)\bigr)\dd t
\]
(both when $w\ne z$, integrating $f'$ along the segment, and when $w=z$), which is continuous because $f'$ is uniformly continuous on $\overline D(z_0,r)$. For fixed $w$, $z\mapsto g(z,w)$ is holomorphic on
$\Omega\setminus\{w\}$ and continuous at $w$, hence holomorphic on $\Omega$ (Volume~5, Theorem~11.20). Writing $\int_\gamma$ as a sum of integrals over parameter intervals, Volume~5, Theorem~11.15(b), shows that
\[
h(z)=\frac1{2\pi i}\int_\gamma g(z,w)\dd w
\]
is holomorphic on $\Omega$. Let $\Omega_1=\{z\notin\gamma^*:n(\gamma,z)=0\}$, an open set, since Volume~7 proves that the winding number is constant on each component of the complement. The hypothesis says exactly that $\Omega\cup\Omega_1=\C$. On $\Omega_1$ let
\[
h_1(z)=\frac1{2\pi i}\int_\gamma\frac{f(w)}{w-z}\dd w,
\]
again holomorphic. For $z\in\Omega\cap\Omega_1$, $h(z)=h_1(z)-f(z)n(\gamma,z)=h_1(z)$. So $h$ and $h_1$ together define an entire function $H$. Volume~7 also proves that the winding number vanishes outside a large disc, so $\Omega_1$
contains $\{|z|>r\}$ for $r=\max_{\gamma^*}|w|$, and there $|H(z)|=|h_1(z)|\le\frac{L(\gamma)\max_{\gamma^*}|f|}{2\pi(|z|-r)}$, where $L(\gamma)$ is the length. So $H$ is bounded and tends to $0$ at infinity; by
Liouville's theorem (Volume~5, Theorem~11.16) it is $0$. For $z\in\Omega\setminus\gamma^*$, $h(z)=0$ is statement (a).

For (b), the set $\gamma^*$ contains no disc: a path of length $L$ is covered by $N$ discs of radius $\frac LN$, of total area $\frac{\pi L^2}N$. So there is $z\in\Omega\setminus\gamma^*$. Apply (a) to
$F(w)=(w-z)f(w)$: $0=F(z)n(\gamma,z)=\frac1{2\pi i}\int_\gamma f(w)\dd w$.
\end{proof}

\begin{beauty}
Dixon's proof never decomposes the region, never triangulates, never deforms a contour. It only says: the function that \emph{ought} to vanish is holomorphic on $\Omega$, and visibly holomorphic and
small on the region where the cycle does not wind; the two regions cover the plane, so Liouville kills it. Topology enters only through the winding number.
\end{beauty}

\begin{corollary}[residue theorem, general form]\label{cor:globalresidue}
Let $\Omega$ be open, $S\subset\Omega$ finite, $f$ holomorphic on $\Omega\setminus S$, and $\gamma$ a cycle in $\Omega\setminus S$ homologous to zero in $\Omega$. Then
\[
\int_\gamma f=2\pi i\sum_{c\in S}n(\gamma,c)\operatorname{Res}_cf.
\]
\end{corollary}
\begin{proof}
Let $P_c(z)=\sum_{k\ge1}a_{-k}(z-c)^{-k}$ be the principal part of the Laurent series of $f$ at $c$ (Volume~5, Theorem~11.22), holomorphic on $\C\setminus\{c\}$. Then $f-\sum_cP_c$ has removable
singularities at the points of $S$, and extends to a holomorphic function on $\Omega$; by Theorem~\ref{thm:globalcauchy}(b) its integral over $\gamma$ is $0$. The series $P_c$ converges uniformly on the
compact set $\gamma^*\subset\C\setminus\{c\}$; the terms with $k\ge2$ have primitives and integrate to $0$ over closed paths, and the term $a_{-1}(z-c)^{-1}$ gives $2\pi i\,n(\gamma,c)a_{-1}$.
\end{proof}

An open set $\Omega$ is \emph{homologically simply connected} if every cycle in $\Omega$ is homologous to zero in $\Omega$: no closed path in $\Omega$ winds around a point outside $\Omega$.

\begin{corollary}\label{cor:hsc}
Let $\Omega$ be a connected, homologically simply connected open set. Then every holomorphic function on $\Omega$ has a primitive, and every holomorphic function without zeros on $\Omega$ has a
holomorphic logarithm, hence holomorphic $n$-th roots for every $n\ge1$.
\end{corollary}
\begin{proof}
Fix $z_0\in\Omega$. Any two points of $\Omega$ are joined by a polygonal path in $\Omega$ (Volume~7, Proposition~3.14), and by Theorem~\ref{thm:globalcauchy}(b) the integral of $f$ along such a path from
$z_0$ to $z$ does not depend on the path; call it $F(z)$. If $D(z,r)\subset\Omega$ and $|h|<r$, then $F(z+h)-F(z)=\int_{[z,z+h]}f$, and $F'(z)=f(z)$ follows exactly as in Volume~5, Theorem~11.9. If $f$
has no zeros, let $L_0$ be a primitive of $\frac{f'}f$. Then $(fe^{-L_0})'=0$, so $f=e^{c+L_0}$ for a constant $c$, and $L=c+L_0$ is a logarithm; $e^{L/n}$ is an $n$-th root.
\end{proof}

\begin{proposition}\label{prop:hsccriterion}
Let $\Omega$ be open. If every point of $\C\setminus\Omega$ lies in a connected unbounded subset of $\C\setminus\Omega$, then $\Omega$ is homologically simply connected. In particular, star-shaped open sets,
the slit plane $\C\setminus\left]-\infty,0\right]$, and the inside of every Jordan curve are homologically simply connected.
\end{proposition}
\begin{proof}
Let $\gamma$ be a closed path in $\Omega$, $a\notin\Omega$, and $E\ni a$ a connected unbounded subset of $\C\setminus\Omega\subset\C\setminus\gamma^*$. The integer-valued continuous function $n(\gamma,\cdot)$ is constant
on $E$, and $E$ contains points of modulus larger than $\max_{\gamma^*}|z|$, where it vanishes (Volume~7). So $n(\gamma,a)=0$.

If $\Omega$ is star-shaped with respect to $c$ and $a\notin\Omega$, the ray $\{a+t(a-c):t\ge0\}$ misses $\Omega$ (if $a+t(a-c)\in\Omega$, the segment from $c$ to it, which contains $a$, would lie in $\Omega$). For the
slit plane, $\C\setminus\Omega$ is itself a connected unbounded set. If $U$ is the inside of a Jordan curve $J$ and $V$ the outside, then $\C\setminus U=J\cup V=\overline V$ by the Jordan curve theorem of Volume~7, the
closure of a connected unbounded set, which is connected and unbounded.
\end{proof}

\begin{think}
The annulus $1<|z|<2$ fails the test, and indeed $\frac1z$ has no primitive there. Can you find a connected open set $\Omega$, with $\C\setminus\Omega$ connected but bounded, that is
homologically simply connected? (Hint: think about what the complement of a closed disc looks like.)
\end{think}

\section{Growth in a strip: Phragm\'en and Lindel\"of}

Lars Edvard Phragm\'en had already made his name in an unusual way: as editor at \emph{Acta Mathematica} in 1889, he found the error in Poincar\'e's prize-winning memoir on the three-body problem,
which had to be pulped and rewritten, and from whose correction chaos theory was born. In 1908, with Ernst Lindel\"of, he extended the maximum principle to unbounded regions. The principle can only
hold with a growth condition, because $\exp(\exp(-is))$ is bounded by $1$ on the edges of the strip $|\operatorname{Re}s|\le\frac\pi2$ and enormous in the middle (Exercise~\ref{ex:plsharp}); but any reasonable growth
is allowed. We write $s=\sigma+it$, as in Volume~6.

\begin{theorem}[Phragm\'en--Lindel\"of in a half-strip]\label{thm:pl}
Let $a<b$, $t_0\in\R$, and $S=\{s:a\le\sigma\le b,\ t\ge t_0\}$. Let $f$ be continuous on $S$ and holomorphic in its interior, with $|f(s)|\le Ce^{Ct}$ on $S$ for some $C$, and $|f|\le M$ on the boundary of $S$.
Then $|f|\le M$ on $S$.
\end{theorem}
\begin{proof}
We may assume $M>0$ (replace $M$ by $M+\eta$ and let $\eta\to0$). Let $\varepsilon>0$, $K=\max(a^2,b^2)$, and $G(s)=f(s)e^{\varepsilon s^2}$. Since $\operatorname{Re}s^2=\sigma^2-t^2$, $|e^{\varepsilon s^2}|\le e^{\varepsilon K}e^{-\varepsilon t^2}$ on $S$. So
$|G|\le Me^{\varepsilon K}$ on the boundary of $S$, and $|G(s)|\le Ce^{Ct+\varepsilon K-\varepsilon t^2}\le Me^{\varepsilon K}$ when $t\ge T$, for some $T$. On the rectangle $\{a\le\sigma\le b,\ t_0\le t\le T'\}$ with $T'\ge T$,
$|G|\le Me^{\varepsilon K}$ on the boundary, hence inside by the maximum principle (Volume~5, Theorem~11.19 and the remark after it). Every $s\in S$ lies in such a rectangle, so
$|f(s)|\le Me^{\varepsilon K}e^{-\varepsilon(\sigma^2-t^2)}$, and letting $\varepsilon\to0$ gives $|f(s)|\le M$.
\end{proof}

\begin{corollary}[convexity]\label{cor:convexity}
Let $S=\{a\le\sigma\le b,\ t\ge1\}$ and $f$ continuous on $S$, holomorphic inside, with $|f(s)|\le Ce^{Ct}$, $|f(a+it)|\le At^\alpha$ and $|f(b+it)|\le Bt^\beta$ for $t\ge1$. Let $\ell$ be the affine function with
$\ell(a)=\alpha$, $\ell(b)=\beta$. Then $|f(\sigma+it)|\le C't^{\ell(\sigma)}$ on $S$, for a constant $C'$.
\end{corollary}
\begin{proof}
Extend $\ell$ to the complex affine function $\ell(s)=\alpha+\kappa(s-a)$, $\kappa=\frac{\beta-\alpha}{b-a}$; then $\operatorname{Re}\ell(s)=\ell(\sigma)$ and $\operatorname{Im}\ell(s)=\kappa t$. On $S$, $-is=t-i\sigma$ has real part
$t\ge1$, so $\operatorname{Log}(-is)$ is holomorphic there, and
\[
\bigl|e^{-\ell(s)\operatorname{Log}(-is)}\bigr|=|s|^{-\ell(\sigma)}e^{\kappa t\arg(-is)},\qquad\arg(-is)=-\arctan\frac\sigma t.
\]
Let $c=\max(|a|,|b|)$ and $L=\max(|\alpha|,|\beta|)$. Since $|t\arctan\frac\sigma t|\le|\sigma|\le c$ and $1\le\frac{|s|}t\le1+c$, this modulus lies between $k^{-1}t^{-\ell(\sigma)}$ and $kt^{-\ell(\sigma)}$, with
$k=e^{|\kappa|c}(1+c)^L$. So $F(s)=f(s)e^{-\ell(s)\operatorname{Log}(-is)}$ satisfies $|F|\le kA$ on $\sigma=a$, $|F|\le kB$ on $\sigma=b$, is bounded on the segment $t=1$ by continuity, and grows at most like
$C''e^{C''t}$. By Theorem~\ref{thm:pl}, $|F|\le M'$ on $S$, and $|f(s)|\le kM't^{\ell(\sigma)}$.
\end{proof}

To apply this to $\zeta$ we need its size on vertical lines, and for that the size of $\Gamma$. Both come from earlier volumes.

\begin{lemma}\label{lem:verticalsizes}
Let $I$ be a compact interval.
\begin{enumerate}
\item[(a)] For $\sigma\in I$ and $|t|\ge1$, $|\Gamma(\sigma+it)|$ lies between two positive constants (depending on $I$) times $|t|^{\sigma-\frac12}e^{-\frac\pi2|t|}$, and
$|\Gamma(\sigma+it)|\sim\sqrt{2\pi}\,|t|^{\sigma-\frac12}e^{-\frac\pi2|t|}$ as $|t|\to\infty$, uniformly in $\sigma\in I$.
\item[(b)] For $|t|\ge1$: $\zeta(\sigma+it)=O(|t|)$ uniformly for $\sigma\ge\frac12$, and $\zeta(\sigma+it)=O(|t|^{\frac32-\sigma})$ uniformly for $\sigma\in I$, $\sigma\le\frac12$.
\end{enumerate}
\end{lemma}
\begin{proof}
(a) Since $\Gamma(\overline s)=\overline{\Gamma(s)}$, take $t\ge1$. For $t\ge T_I=\max(1,\max_I|\sigma|)$, $|\arg s|\le\frac{3\pi}4$, and Stirling's formula (Volume~6, Theorem~6.1(c)) gives
$\log|\Gamma(s)|=\operatorname{Re}\bigl((s-\frac12)\log s\bigr)-\sigma+\frac12\log2\pi+O(\frac1t)$. Here $\log|s|=\log t+O(t^{-2})$ and $\arg s=\frac\pi2-\arctan\frac\sigma t=\frac\pi2-\frac\sigma t+O(t^{-3})$, so
$\operatorname{Re}\bigl((s-\frac12)\log s\bigr)=(\sigma-\frac12)\log|s|-t\arg s=(\sigma-\frac12)\log t-\frac\pi2t+\sigma+O(\frac1t)$. Hence $\log|\Gamma(s)|=(\sigma-\frac12)\log t-\frac\pi2t+\frac12\log2\pi+O(\frac1t)$,
uniformly, which is the asymptotic statement. On the compact set $I\times[1,T_I]$ the ratio of $|\Gamma|$ to the model is continuous and positive, so bounded above and below.

(b) In the proof of Volume~5, Proposition~12.6, $\zeta(s)-\frac1{s-1}=\sum_n\phi_n(s)$ with $|\phi_n(s)|\le|s|n^{-\sigma-1}$. For $\frac12\le\sigma\le2$ and $|t|\ge1$ this gives $|\zeta(s)|\le1+(2+|t|)\zeta(\frac32)=O(|t|)$, and
for $\sigma\ge2$, $|\zeta(s)|\le\zeta(2)$. For $\sigma\le\frac12$ put $w=1-s$, so $\operatorname{Re}w\ge\frac12$, and use Volume~5, Corollary~12.9(d):
$\zeta(s)=2(2\pi)^{-w}\cos(\frac{\pi w}2)\Gamma(w)\zeta(w)$. Here $|(2\pi)^{-w}|=(2\pi)^{\sigma-1}$ is bounded for $\sigma\in I$, $|\cos\frac{\pi w}2|\le\cosh\frac{\pi t}2\le e^{\frac\pi2|t|}$, $|\Gamma(w)|=O(|t|^{\frac12-\sigma}e^{-\frac\pi2|t|})$
by (a), and $\zeta(w)=O(|t|)$. Multiply.
\end{proof}

\begin{theorem}[the convexity bound]\label{thm:convexzeta}
For every $\varepsilon>0$, $\zeta(\frac12+it)=O_\varepsilon(|t|^{\frac14+\varepsilon})$ for $|t|\ge1$.
\end{theorem}
\begin{proof}
By $\zeta(\overline s)=\overline{\zeta(s)}$, take $t\ge1$. Let $0<\varepsilon<\frac12$ and apply Corollary~\ref{cor:convexity} on $S=\{-\varepsilon\le\sigma\le1+\varepsilon,\ t\ge1\}$, where $\zeta$ is holomorphic (its only pole is
at $1$) and of polynomial growth by Lemma~\ref{lem:verticalsizes}(b). On $\sigma=1+\varepsilon$, $|\zeta(s)|\le\sum_nn^{-1-\varepsilon}=\zeta(1+\varepsilon)$: take $\alpha=0$. On $\sigma=-\varepsilon$, the functional equation with
$w=1+\varepsilon-it$, as in the lemma, gives $|\zeta(s)|\le2(2\pi)^{-1-\varepsilon}e^{\frac\pi2t}|\Gamma(w)|\zeta(1+\varepsilon)=O(t^{\frac12+\varepsilon})$: take $\beta=\frac12+\varepsilon$. At $\sigma=\frac12$, which is at fraction
$\frac{1/2+\varepsilon}{1+2\varepsilon}=\frac12$ of the way from $1+\varepsilon$ to $-\varepsilon$, the interpolated exponent is $\frac12(\frac12+\varepsilon)=\frac14+\frac\varepsilon2$.
\end{proof}

\begin{remark}
The \emph{Lindel\"of hypothesis} says that the exponent $\frac14$ can be replaced by $0$: $\zeta(\frac12+it)=O(|t|^\varepsilon)$ for every $\varepsilon>0$. It follows from the Riemann hypothesis. Weyl's
exponential sum method, in the hands of Hardy and Littlewood, gave $\frac16$ around 1921 (see Titchmarsh's book, Chapter~14), and a century of work brought it down to $\frac{13}{84}\approx0.155$ (Bourgain, 2017). Everything beyond
$\frac14$ requires arithmetic: the convexity bound is what complex analysis gives for free.
\end{remark}

\section{Entire functions of finite order}

A polynomial is determined by its roots up to a constant factor. In 1876 Weierstrass showed that any sequence of points tending to infinity is the zero set of an entire function, built as an
infinite product; the factors need small exponential corrections to make the product converge. In 1893 Hadamard showed that when the function does not grow too fast, the corrections can be
taken of bounded degree, and what remains is the exponential of a polynomial. Volume~6 (Theorem~6.7) proved this for Riemann's $\xi$; here is the general theorem.

An entire function $f$ has \emph{finite order} if $|f(z)|\le\exp(|z|^\lambda)$ for some $\lambda$ and all large $|z|$; its \emph{order} $\rho$ is the infimum of such $\lambda$. Polynomials have order $0$; $e^z$, $\sin z$
and $\xi$ have order $1$; $e^{z^2}$ has order $2$; $\exp(e^z)$ has infinite order. With $M(r)=\max_{|z|=r}|f|$, the order is $\limsup_{r\to\infty}\frac{\log\log M(r)}{\log r}$.

For $p\ge0$, the \emph{elementary factor} is $E_p(z)=(1-z)\exp\bigl(z+\frac{z^2}2+\dots+\frac{z^p}p\bigr)$; $E_0(z)=1-z$. Its exponent is the beginning of the series of $-\log(1-z)$, so $E_p(z)$ is close to $1$
for small $z$.

\begin{lemma}\label{lem:products}
\begin{enumerate}
\item[(a)] For $|z|\le1$, $|1-E_p(z)|\le|z|^{p+1}$.
\item[(b)] Let $u_n$ be holomorphic on an open set $\Omega$, with $\sum_n|u_n|$ converging uniformly on every compact subset. Then $\prod_{n\le N}(1+u_n)$ converges uniformly on compact sets to a holomorphic
$P$, and $P(z)=0$ exactly when some $1+u_n(z)=0$, the multiplicity of the zero being the sum of the multiplicities in the factors.
\end{enumerate}
\end{lemma}
\begin{proof}
(a) $E_p'(z)=-z^p\exp\bigl(z+\dots+\frac{z^p}p\bigr)$, so $-E_p'$ has a power series with non-negative coefficients starting at $z^p$. Since $E_p(0)=1$, $1-E_p(z)=z^{p+1}\varphi(z)$ with $\varphi$ a power series
with non-negative coefficients, and $\varphi(1)=1-E_p(1)=1$. For $|z|\le1$, $|1-E_p(z)|\le|z|^{p+1}\varphi(|z|)\le|z|^{p+1}\varphi(1)$.

(b) Let $K$ be compact, and $N_0$ such that $\sum_{n>N_0}|u_n|\le\frac12$ on $K$. For $|u|\le\frac12$, the series $\log(1+u)=\sum_{k\ge1}\frac{(-1)^{k+1}u^k}k$ satisfies $|\log(1+u)|\le\frac{|u|}{1-|u|}\le2|u|$. So
$\sum_{n>N_0}\log(1+u_n)$ converges uniformly on $K$ to a bounded $\Sigma$, and $\prod_{N_0<n\le N}(1+u_n)=\exp\sum_{N_0<n\le N}\log(1+u_n)\to e^\Sigma$ uniformly. Hence on $K$, $P=\prod_{n\le N_0}(1+u_n)\cdot e^\Sigma$, where
$e^\Sigma$ never vanishes. Covering $\Omega$ by discs with compact closures, $P$ is holomorphic (Volume~5, Theorem~11.15(a)) with the stated zeros.
\end{proof}

\begin{lemma}\label{lem:zerocount}
Let $f$ be entire of finite order $\rho$, not identically zero, with non-zero zeros $a_1,a_2,\dots$ (repeated with multiplicity). Then $\sum_n|a_n|^{-\lambda}<\infty$ for every $\lambda>\rho$.
\end{lemma}
\begin{proof}
Dividing by a power of $z$ does not increase the order, so assume $f(0)\ne0$. Choose $\eta>0$ with $\rho+\eta<\lambda$. By Volume~6, Corollary~6.4, with radii $r$ and $2r$, the number of zeros in $|z|\le r$ is
$n(r)\le\frac{\log M(2r)-\log|f(0)|}{\log2}=O(r^{\rho+\eta})$. The zeros with $2^k\le|a_n|<2^{k+1}$ contribute at most $n(2^{k+1})2^{-k\lambda}=O(2^{(k+1)(\rho+\eta)-k\lambda})$, the terms of a convergent geometric series, and
there are finitely many zeros in $|z|<1$.
\end{proof}

\begin{theorem}[Hadamard factorisation]\label{thm:hadamardgen}
Let $f$ be entire of finite order $\rho$, with a zero of order $m\ge0$ at $0$ and other zeros $a_1,a_2,\dots$, repeated with multiplicity, and let $p=\lfloor\rho\rfloor$. Then
\[
f(z)=z^me^{Q(z)}\prod_nE_p\Bigl(\frac z{a_n}\Bigr),
\]
where the product converges uniformly on compact sets and $Q$ is a polynomial of degree at most $p$.
\end{theorem}
\begin{proof}
\emph{The product.} For $|z|\le r$ and $|a_n|\ge r$, $|1-E_p(\frac z{a_n})|\le r^{p+1}|a_n|^{-p-1}$ by Lemma~\ref{lem:products}(a), and $\sum_n|a_n|^{-p-1}<\infty$ by Lemma~\ref{lem:zerocount} since $p+1>\rho$. By
Lemma~\ref{lem:products}(b), the product $P$ is entire and vanishes exactly at the $a_n$, with the right multiplicities ($E_p(\frac z a)$ has a single simple zero, at $a$). So $\frac f{z^mP}$ extends to an entire
function without zeros (Volume~5, Theorem~11.20), which has a logarithm $Q$ by Corollary~\ref{cor:hsc}: $f=z^me^QP$.

\emph{The logarithmic derivative.} The identity $\frac1{z-a}+\sum_{k=0}^{p-1}\frac{z^k}{a^{k+1}}=\frac{z^p}{a^p(z-a)}$ and the formula $\frac{E_p'}{E_p}(u)=\frac1{u-1}+\sum_{k=0}^{p-1}u^k$ give
\[
\frac{f'}f(z)=\frac mz+Q'(z)+\sum_n\frac{z^p}{a_n^p(z-a_n)}
\]
away from the zeros: the partial products converge locally uniformly, so do their derivatives (Volume~5, Theorem~11.15(a)), hence so do their logarithmic derivatives on compact sets where $P\ne0$.

\emph{Killing $Q^{(p+1)}$.} Fix $R\ge1$, and let $g_R(z)=\frac{f(z)}{z^m\prod_{|a_n|\le2R}(1-z/a_n)}$, an entire function without zeros in $|z|\le2R$, with $g_R(0)=c$ independent of $R$ ($c$ is the first non-zero
Taylor coefficient of $f$). On $|z|=4R$ each factor satisfies $|1-\frac z{a_n}|\ge\frac{|z|}{|a_n|}-1\ge1$ and $|z^m|\ge1$, so $|g_R|\le M(4R)$ there, and on $|z|\le4R$ by the maximum principle. On the disc $|z|<2R$, $g_R$ has
a holomorphic logarithm $L_R$ (Corollary~\ref{cor:hsc}), and $h=L_R-L_R(0)$ satisfies $h(0)=0$ and $\operatorname{Re}h=\log|g_R|-\log|c|\le A_R=\log M(4R)-\log|c|$. Borel--Carath\'eodory (Volume~6, Lemma~6.5) on the
radius $\frac32R$ gives, for $|z|\le R$,
\[
|L_R^{(p+1)}(z)|\le\frac{2A_R\,(p+1)!\,\frac32R}{(R/2)^{p+2}}=O\Bigl(\frac{A_R}{R^{p+1}}\Bigr)\longrightarrow0\quad(R\to\infty),
\]
because $A_R=O(R^{\rho+\eta})$ with $\rho+\eta<p+1$. On the other hand $L_R'=\frac{f'}f-\frac mz-\sum_{|a_n|\le2R}\frac1{z-a_n}$, and by the formula above (whose series, for $|z|\le R$ and $|a_n|>2R$, has terms at most $2R^p|a_n|^{-p-1}$ and may be differentiated term by term), differentiating $p$ more
times kills the polynomial terms $\sum_{k<p}\frac{z^k}{a_n^{k+1}}$:
\[
L_R^{(p+1)}(z)=Q^{(p+1)}(z)+(-1)^pp!\sum_{|a_n|>2R}\frac1{(z-a_n)^{p+1}}\qquad(|z|<2R).
\]
For $|z|\le R$ the last sum is bounded by $\sum_{|a_n|>2R}(\frac{|a_n|}2)^{-p-1}$, a tail of a convergent series. Letting $R\to\infty$, $Q^{(p+1)}(z)=0$. So $Q$ is a polynomial of degree at most $p$.
\end{proof}

\begin{beauty}
Jensen's formula says a slowly growing function has few zeros; Borel--Carath\'eodory says a holomorphic function whose real part is small has small derivatives. Hadamard's theorem is the meeting
of the two: the zeros are few enough for a product with corrections of degree $p$ to converge, and what is left over, the polynomial $Q$, is pinned down by a bound on its real part alone.
\end{beauty}

\section{The Mellin transform and harmonic sums}\label{sec:mellin}

The Mellin transform is the Fourier transform of the multiplicative group $]0,\infty[$, whose Haar measure is $\frac{\dd x}x$ (Chapter~21). Robert Hjalmar Mellin, a Finnish student of Mittag-Leffler,
made it systematic in the 1890s; Riemann had already used it in 1859 when he wrote $\Gamma(\frac s2)\pi^{-s/2}\zeta(s)$ as an integral of the theta function (Volume~5, Lemma~12.7). A century later Philippe
Flajolet made it the workhorse of the analysis of algorithms, because it turns a sum of dilations, the shape of a cost that repeats at scales $2,4,8,\dots$, into a product.

For $f:]0,\infty[\to\C$ continuous, its \emph{Mellin transform} is
\[
\mathcal Mf(s)=\int_0^\infty f(x)x^{s-1}\dd x,
\]
defined where the integral converges absolutely. If $\int_0^\infty|f(x)|x^{\sigma-1}\dd x<\infty$ for $\sigma$ in an open interval $]\sigma_1,\sigma_2[$, then $\mathcal Mf$ is holomorphic on the strip $\sigma_1<\sigma<\sigma_2$
(Volume~5, Theorem~11.15(b): for $\sigma'\le\sigma\le\sigma''$ inside the strip, $|f(x)x^{s-1}|\le|f(x)|(x^{\sigma'-1}+x^{\sigma''-1})$). Examples: $\mathcal M(e^{-x})=\Gamma(s)$ for $\sigma>0$; integrating by parts,
for $-1<\sigma<0$, where the boundary terms vanish,
\begin{equation}\label{eq:mellin1me}
\mathcal M(1-e^{-x})(s)=-\frac1s\int_0^\infty e^{-x}x^s\dd x=-\Gamma(s).
\end{equation}

\begin{theorem}[Mellin inversion]\label{thm:mellininv}
Let $f:]0,\infty[\to\C$ be continuous, and $c$ real with $\int_0^\infty|f(x)|x^{c-1}\dd x<\infty$ and $\int_{-\infty}^\infty|\mathcal Mf(c+i\tau)|\dd\tau<\infty$. Then for every $x>0$,
\[
f(x)=\frac1{2\pi i}\int_{c-i\infty}^{c+i\infty}\mathcal Mf(s)x^{-s}\dd s.
\]
\end{theorem}
\begin{proof}
Let $F(u)=f(e^u)e^{cu}$. The change of variables $x=e^u$ gives $\int|F|=\int_0^\infty|f(x)|x^{c-1}\dd x<\infty$ and $\hat F(\xi)=\int F(u)e^{-2\pi iu\xi}\dd u=\mathcal Mf(c-2\pi i\xi)$, which is integrable. $F$ is continuous, so
Fourier inversion (Theorem~\ref{thm:ftinv}) gives $F(u)=\int\hat F(\xi)e^{2\pi iu\xi}\dd\xi$. Substitute $\tau=-2\pi\xi$ and $u=\log x$: $f(x)x^c=\frac1{2\pi}\int\mathcal Mf(c+i\tau)x^{-i\tau}\dd\tau$, which is the
formula, with $s=c+i\tau$, $\dd s=i\dd\tau$.
\end{proof}

A \emph{harmonic sum} is $G(x)=\sum_k\lambda_kg(\mu_kx)$, with amplitudes $\lambda_k$ and frequencies $\mu_k>0$. If $\sum_k|\lambda_k|\mu_k^{-\sigma}<\infty$ and $\int|g|x^{\sigma-1}<\infty$, the substitution $y=\mu_kx$ and
Fubini give the separation of variables
\begin{equation}\label{eq:harmonicsum}
\mathcal MG(s)=\Lambda(s)\,\mathcal Mg(s),\qquad\Lambda(s)=\sum_k\lambda_k\mu_k^{-s}.
\end{equation}
The Dirichlet series $\Lambda$ carries the arithmetic of the frequencies, $\mathcal Mg$ the shape of the base function. Asymptotics of $G$ then come from moving the line of integration in
Theorem~\ref{thm:mellininv} across poles: as $x\to0$ to the left, as $x\to\infty$ to the right, each pole at $s_0$ contributing a term $x^{-s_0}$ (times powers of $\log x$ at multiple poles).

\begin{example}[a theta sum]\label{ex:thetamellin}
Let $G(x)=\sum_{k\ge1}e^{-k^2x}$. Then for every $N$, as $x\to0^+$, $G(x)=\frac12\sqrt{\frac\pi x}-\frac12+O(x^N)$.

Here $g(x)=e^{-x}$ and $\mu_k=k^2$, so $\mathcal MG(s)=\Gamma(s)\zeta(2s)$ for $\sigma>\frac12$. Since $G(x)\le\int_0^\infty e^{-u^2x}\dd u=\frac12\sqrt{\pi/x}$ and $G(x)\le\frac{e^{-x}}{1-e^{-x}}$, the hypotheses of
Theorem~\ref{thm:mellininv} hold at $c=1$ (by Lemma~\ref{lem:verticalsizes}, $\Gamma$ decays exponentially on vertical lines while $\zeta(2s)$ grows at most polynomially). Integrate $\Gamma(s)\zeta(2s)x^{-s}$ around the
rectangle with vertical sides $\sigma=1$ and $\sigma=-N-\frac12$ and horizontal sides $t=\pm T$. The horizontal sides contribute $O(e^{-T})$ by the same lemma, which tends to $0$. The poles inside are at
$s=\frac12$, where $\zeta(2s)=\frac1{2(s-\frac12)}+O(1)$, with residue $\frac12\Gamma(\frac12)x^{-1/2}=\frac12\sqrt{\pi/x}$; and at $s=0$, residue $\zeta(0)=-\frac12$ (Volume~5, Corollary~12.9). The poles of $\Gamma$ at
$-1,-2,\dots$ are cancelled by the trivial zeros $\zeta(-2),\zeta(-4),\dots$. So $G(x)$ equals these two residues plus $\frac1{2\pi i}\int_{(-N-\frac12)}\Gamma(s)\zeta(2s)x^{-s}\dd s$, which is $O(x^{N+\frac12})$.
\end{example}

The theta functional equation (Volume~5, Chapter~12) shows that the error is in fact $\sqrt{\pi/x}\sum_{k\ge1}e^{-\pi^2k^2/x}$, exponentially small: the Mellin method found every power of $x$ there is, namely none.
Our main example is from computer science.

\emph{Probabilistic counting.} How many distinct visitors did a website have today? Storing them all is expensive. In 1985 Flajolet and Nigel Martin suggested hashing each visitor to a random bit
string and remembering only the largest number of leading zeros seen. A visitor's \emph{level} $L$ satisfies $\Pr(L\ge k)=2^{-k}$, and with $n$ visitors the maximum level is about $\log_2n$; HyperLogLog
(Flajolet, Fusy, Gandouet and Meunier, 2007), refined from this idea, is now used in many large databases to count distinct items with a few kilobytes of memory. To correct the estimate one needs the
mean of the maximum precisely. It is easiest with a Poisson number of items: let $N$ be Poisson with mean $x$, give the $N$ items independent levels, and let $H$ be the maximum level ($0$ if $N=0$).
Then
\[
\Pr(H<k)=\sum_{n\ge0}e^{-x}\frac{x^n}{n!}(1-2^{-k})^n=e^{-x2^{-k}},\qquad\mathbb E[H]=\sum_{k\ge1}\Pr(H\ge k)=\sum_{k\ge1}\bigl(1-e^{-x2^{-k}}\bigr).
\]

\begin{theorem}[the maximum of geometric levels]\label{thm:maxgeom}
As $x\to\infty$, for every $d>0$,
\[
\mathbb E[H]=\log_2x+\frac\gamma{\ln2}-\frac12+P(\log_2x)+O(x^{-d}),\qquad P(u)=-\frac1{\ln2}\sum_{k\ne0}\Gamma\Bigl(\frac{2\pi ik}{\ln2}\Bigr)e^{-2\pi iku},
\]
where $P$ is a continuous $1$-periodic function with $|P|<1.6\cdot10^{-6}$.
\end{theorem}
\begin{proof}
$G(x)=\mathbb E[H]$ is the harmonic sum with $g(x)=1-e^{-x}$, $\lambda_k=1$ and $\mu_k=2^{-k}$, so by \eqref{eq:mellin1me} and \eqref{eq:harmonicsum}, for $-1<\sigma<0$,
\[
\mathcal MG(s)=-\Gamma(s)\sum_{k\ge1}2^{ks}=-\Gamma(s)\frac{2^s}{1-2^s}.
\]
On the line $\sigma=c=-\frac12$, $|\frac{2^s}{1-2^s}|\le\frac{2^c}{1-2^c}$ and $\Gamma$ decays exponentially, so Theorem~\ref{thm:mellininv} applies. Move the line to $\sigma=d>0$. The function $\mathcal MG$ is meromorphic, with poles
only where $1-2^s=0$, at $\chi_k=\frac{2\pi ik}{\ln2}$ (double at $\chi_0=0$, where $\Gamma$ also has a pole). Use the rectangles $c\le\sigma\le d$, $|t|\le T_j=\frac{(2j+1)\pi}{\ln2}$: on their horizontal sides $2^{iT_j}=-1$, so
$|1-2^s|=1+2^\sigma\ge1$, and the exponential decay of $\Gamma$ (Lemma~\ref{lem:verticalsizes}) makes these sides contribute nothing in the limit. On $\sigma=d$, $|\frac{2^s}{1-2^s}|\le\frac{2^d}{2^d-1}$. Going around the
rectangle counterclockwise (up along $\sigma=d$, down along $\sigma=c$), the residue theorem gives
\[
G(x)=-\sum_{k\in\Z}\operatorname{Res}_{s=\chi_k}\bigl(\mathcal MG(s)x^{-s}\bigr)+\frac1{2\pi i}\int_{(d)}\mathcal MG(s)x^{-s}\dd s,
\]
and the last integral is $O(x^{-d})$ since $|x^{-s}|=x^{-d}$ there. At $\chi_k$ with $k\ne0$, $1-2^s$ has a simple zero with derivative $-\ln2$, so the residue is $\frac{\Gamma(\chi_k)}{\ln2}x^{-\chi_k}$, and
$x^{-\chi_k}=e^{-2\pi ik\log_2x}$: these give $P(\log_2x)$. At $0$, with $\Gamma(s)=\frac1s-\gamma+O(s)$ (as $\Gamma'(1)=-\gamma$, Volume~5, Exercise~7.7), $\frac{2^s}{1-2^s}=-\frac1{s\ln2}-\frac12+O(s)$ and $x^{-s}=1-s\ln x+O(s^2)$,
\[
\begin{aligned}\mathcal MG(s)x^{-s}&=\Bigl(\frac1s-\gamma\Bigr)\Bigl(\frac1{s\ln2}+\frac12\Bigr)(1-s\ln x)+O(1)\\&=\frac1{s^2\ln2}+\frac1s\Bigl(\frac12-\frac\gamma{\ln2}-\log_2x\Bigr)+O(1),\end{aligned}
\]
whose residue, with the minus sign, gives the main terms. Finally, the reflection formula (Volume~5, Theorem~7.13) and $\Gamma(1-iy)=-iy\Gamma(-iy)$ give $|\Gamma(iy)|^2=\frac\pi{y\sinh\pi y}$. With $y_k=\frac{2\pi k}{\ln2}$, so
$y_1\approx9.06$, $|P|\le\frac2{\ln2}\sum_{k\ge1}\bigl(\frac\pi{y_k\sinh\pi y_k}\bigr)^{1/2}<1.6\cdot10^{-6}$; the series converges so fast that $P$ is continuous (indeed analytic).
\end{proof}

\begin{beauty}
The mean of the maximum is not quite $\log_2x+0.3327$: it wobbles periodically in $\log_2x$, by a millionth. The wobble is real, it comes from the poles of $\frac1{1-2^s}$ spaced along the imaginary
axis, and it is invisible in any simulation, while the Mellin transform sees it at once. Such fluctuations are typical of algorithms that work by halving: tries, skip lists, radix sort
and approximate counting all have them, and Flajolet, Gourdon and Dumas's survey (see the sources) computes many of them.
\end{beauty}

\section{Exercises}

\begin{exercise}[$\star\star$]\label{ex:annuluscauchy}
\emph{Laurent series from the homology form.} Let $f$ be holomorphic on the annulus $\Omega=\{r<|z|<R\}$, and $r<r_1<r_2<R$. Let $C(\rho)$ be the circle $|z|=\rho$, counterclockwise. (a) Show that the cycle
$C(r_2)-C(r_1)$ is homologous to zero in $\Omega$, and that for $r_1<|z|<r_2$,
\[
f(z)=\frac1{2\pi i}\int_{C(r_2)}\frac{f(w)}{w-z}\dd w-\frac1{2\pi i}\int_{C(r_1)}\frac{f(w)}{w-z}\dd w.
\]
(b) Expand both kernels in geometric series to obtain the Laurent expansion $f(z)=\sum_{k\in\Z}a_kz^k$ on the annulus.
\end{exercise}

\begin{exercise}[$\star$]\label{ex:plsharp}
\emph{Growth conditions are necessary.} Let $f(s)=\exp(\exp(-is))$ on the strip $|\sigma|\le\frac\pi2$. (a) Show that $|f|=1$ on both edges and $|f(it)|=\exp(e^t)$. (b) Which hypothesis of
Theorem~\ref{thm:pl} fails?
\end{exercise}

\begin{exercise}[$\star\star$]\label{ex:sinprod}
\emph{Euler's sine product from Hadamard.} (a) Show that $\sin\pi z$ has order exactly $1$. (b) Deduce from Theorem~\ref{thm:hadamardgen} that $\sin\pi z=\pi z\prod_{n\ge1}\bigl(1-\frac{z^2}{n^2}\bigr)$. (c) Compare the
coefficients of $z^3$ and recover $\zeta(2)=\frac{\pi^2}6$, as Euler did in 1735, this time with a proof that the product is right.
\end{exercise}

\begin{exercise}[$\star\star$]\label{ex:nonintorder}
(a) Show that an entire function of finite order with only finitely many zeros has an integer order. (b) Deduce that $\cos\sqrt z=\sum_{n\ge0}\frac{(-1)^nz^n}{(2n)!}$, which has order $\frac12$, has infinitely
many zeros, without finding them.
\end{exercise}

\begin{exercise}[$\star\star$]\label{ex:lambert}
\emph{Divisors and a Lambert series.} Show that, as $x\to0^+$,
\[
\sum_{k\ge1}\frac1{e^{kx}-1}=\sum_{n\ge1}d(n)e^{-nx}=\frac{\log\frac1x+\gamma}x+\frac14-\frac x{144}+O(x^{3/2}),
\]
where $d(n)$ is the number of divisors of $n$.
\end{exercise}

\section{Hints}

\begin{hint}{ex:annuluscauchy}
Compute the winding number of the cycle around points with $|a|\le r$, with $|a|\ge R$, and around $z$.
\end{hint}
\begin{hint}{ex:plsharp}
Compute the real part of $\exp(-is)$ for $s=\sigma+it$.
\end{hint}
\begin{hint}{ex:sinprod}
Look at $\sin\pi z$ on the imaginary axis. With $p=1$, pair the factors for $n$ and $-n$, and use that $\sin$ is odd.
\end{hint}
\begin{hint}{ex:nonintorder}
By Hadamard, such a function is a polynomial times $e^{Q}$. Choose a direction in which the leading term of $Q$ is real and positive.
\end{hint}
\begin{hint}{ex:lambert}
The Mellin transform is $\Gamma(s)\zeta(s)^2$. Move the line from $\sigma=3$ to $\sigma=-\frac32$ and collect the poles at $1$, $0$ and $-1$.
\end{hint}

\section{Solutions}

\begin{solution}{ex:annuluscauchy}
\idea The two circles together do not wind around the hole, nor around anything outside, so Theorem~\ref{thm:globalcauchy} applies to the annulus even though it is not star-shaped.
\writeup (a) For $|a|\le r$, both circles wind once around $a$, so the cycle has winding number $1-1=0$; for $|a|\ge R$, both wind $0$ times. So the cycle is homologous to zero in $\Omega$. For
$r_1<|z|<r_2$, $n(C(r_2),z)=1$ and $n(C(r_1),z)=0$, and Theorem~\ref{thm:globalcauchy}(a) gives the formula. (b) For $|w|=r_2$, $\frac1{w-z}=\sum_{k\ge0}\frac{z^k}{w^{k+1}}$, and for $|w|=r_1$,
$-\frac1{w-z}=\sum_{k\ge0}\frac{w^k}{z^{k+1}}$, both uniformly in $w$ on the circle. Integrating term by term, $f(z)=\sum_{k\ge0}a_kz^k+\sum_{k\ge1}a_{-k}z^{-k}$ with $a_k=\frac1{2\pi i}\int_{C(r_2)}f(w)w^{-k-1}\dd w$ and
$a_{-k}=\frac1{2\pi i}\int_{C(r_1)}f(w)w^{k-1}\dd w$. By Volume~5, Lemma~11.21, these do not depend on the radii, so the expansion holds on the whole annulus.
\end{solution}

\begin{solution}{ex:plsharp}
\idea The function is built to be harmless on the edges and explosive in the middle, so it exhibits the only loophole the theorem leaves open.
\writeup (a) $\exp(-is)=\exp(t-i\sigma)=e^t(\cos\sigma-i\sin\sigma)$, so $|f(s)|=\exp(e^t\cos\sigma)$. On $\sigma=\pm\frac\pi2$ this is $1$; at $\sigma=0$ it is $\exp(e^t)$. (b) Everything holds (continuity, holomorphy,
$|f|\le1$ on the edges and $|f|\le e^e$ on any bottom segment) except the growth condition $|f(s)|\le Ce^{Ct}$: $f$ grows doubly exponentially. So some growth condition is indispensable. (In fact
Phragm\'en and Lindel\"of showed that $|f|\le\exp(Ce^{\kappa t})$ with $\kappa<1$ suffices on this strip of width $\pi$, so this example is exactly at the border.)
\end{solution}

\begin{solution}{ex:sinprod}
\idea Hadamard reduces the identity to two unknown constants, and symmetry and the value at $0$ determine them.
\writeup (a) $|\sin\pi z|\le\cosh(\pi\operatorname{Im}z)\le e^{\pi|z|}$, so the order is at most $1$; and $|\sin(\pi iy)|=\sinh\pi y\ge\frac14e^{\pi y}$ for large $y$, so it is at least $1$. (b) The zeros are the integers, all simple.
By Theorem~\ref{thm:hadamardgen} with $p=1$, $\sin\pi z=ze^{A+Bz}\prod_{n\ne0}(1-\frac zn)e^{z/n}$. The product converges absolutely, so we may group $n$ with $-n$: $(1-\frac zn)e^{z/n}(1+\frac zn)e^{-z/n}=1-\frac{z^2}{n^2}$. Hence
$\sin\pi z=ze^{A+Bz}\Pi(z)$ with $\Pi$ even. Replacing $z$ by $-z$ and using oddness gives $e^{Bz}=e^{-Bz}$ wherever $\Pi(z)\ne0$, so $B=0$ (differentiate at $0$). Dividing by $z$ and letting $z\to0$, $e^A=\pi$.
(c) The product converges uniformly near $0$, so its Taylor coefficients are limits of those of the partial products: the coefficient of $z^2$ in $\Pi$ is $-\sum_n\frac1{n^2}$. Comparing with
$\sin\pi z=\pi z-\frac{\pi^3z^3}6+\dots$ gives $\sum_n\frac1{n^2}=\frac{\pi^2}6$.
\end{solution}

\begin{solution}{ex:nonintorder}
\idea Without zeros to account for, the growth of $f$ is all in $e^Q$, and $e^{Q}$ grows like $\exp(|z|^{\deg Q})$ in some direction.
\writeup (a) With finitely many zeros, Theorem~\ref{thm:hadamardgen} (where the product is finite) gives $f=Pe^Q$ with $P$ a polynomial and $Q$ a polynomial of degree $d$. If $d=0$, $f$ is a polynomial of order
$0$. If $d\ge1$, write $Q(z)=cz^d+\dots$ with $c\ne0$, and choose $\theta$ with $ce^{id\theta}=|c|$. Then $\operatorname{Re}Q(re^{i\theta})\ge\frac{|c|}2r^d$ and $|P(re^{i\theta})|\ge1$ for large $r$, so $M(r)\ge\exp(\frac{|c|}2r^d)$ and the order is
at least $d$; it is at most $d$ since $|Pe^Q|\le\exp(C|z|^d)$ for large $|z|$. So the order is the integer $d$. (b) $|\cos\sqrt z|\le\cosh|\sqrt z|\le e^{|z|^{1/2}}$ (the series does not depend on the choice of
square root), and $\cos\sqrt{-r}=\cosh\sqrt r\ge\frac12e^{\sqrt r}$, so the order is $\frac12$, not an integer. By (a), the zeros are infinitely many: they are $\bigl((k+\frac12)\pi\bigr)^2$, $k\ge0$.
\end{solution}

\begin{solution}{ex:lambert}
\idea A harmonic sum with base function $e^{-x}$ and frequencies $n$ weighted by $d(n)$: its Mellin transform is $\Gamma(s)$ times the Dirichlet series of $d$, which is $\zeta(s)^2$. Each pole gives one term.
\writeup Expanding $\frac1{e^{kx}-1}=\sum_{m\ge1}e^{-kmx}$ and grouping $n=km$ gives the first equality. For $\sigma>1$, $\sum_nd(n)n^{-s}=\zeta(s)^2$ (Volume~5, Proposition~12.2), so $\mathcal MG(s)=\Gamma(s)\zeta(s)^2$. As
$G(x)=O(x^{-2})$ near $0$ (since $d(n)\le n$) and $G$ decays exponentially at infinity, Theorem~\ref{thm:mellininv} applies at $c=3$, using Lemma~\ref{lem:verticalsizes}. Shift to $\sigma=-\frac32$ as in
Example~\ref{ex:thetamellin}; the new integral is $O(x^{3/2})$. Residues: at $s=1$, with $\Gamma(s)=1-\gamma(s-1)+\dots$ (Volume~5, Exercise~7.7), $\zeta(s)^2=\frac1{(s-1)^2}+\frac{2\gamma}{s-1}+\dots$ (Volume~5,
Proposition~10.6) and $x^{-s}=\frac1x\bigl(1-(s-1)\log x+\dots\bigr)$, the coefficient of $\frac1{s-1}$ is $\frac1x(-\gamma-\log x+2\gamma)=\frac{\log\frac1x+\gamma}x$. At $s=0$: $\operatorname{Res}\Gamma=1$ and $\zeta(0)^2=\frac14$. At $s=-1$:
$\operatorname{Res}\Gamma=-1$ and $\zeta(-1)^2=\frac1{144}$ (Volume~5, Theorem~12.10), giving $-\frac x{144}$.
\end{solution}

\section*{Sources and further reading}
\addcontentsline{toc}{section}{Sources and further reading}

J.~D.~Dixon, ``A brief proof of Cauchy's integral theorem'', \emph{Proc.\ Amer.\ Math.\ Soc.} 29 (1971), 625--626. E.~Phragm\'en and E.~Lindel\"of, ``Sur une extension d'un principe classique de
l'analyse'', \emph{Acta Math.} 31 (1908), 381--406. Phragm\'en's discovery of the error in Poincar\'e's prize memoir is told in J.~Barrow-Green, \emph{Poincar\'e and the Three
Body Problem} (AMS and LMS, 1997). J.~Hadamard, ``\'Etude sur les propri\'et\'es des fonctions enti\`eres et en particulier d'une fonction consid\'er\'ee par Riemann'', \emph{J.\ Math.\
Pures Appl.} (4) 9 (1893), 171--215. For the Lindel\"of hypothesis and the subconvexity record: J.~Bourgain, ``Decoupling, exponential sums and the Riemann zeta function'', \emph{J.\ Amer.\ Math.\
Soc.} 30 (2017), 205--224; and E.~C.~Titchmarsh, \emph{The Theory of the Riemann Zeta-Function}, 2nd ed. (Oxford, 1986), Chapter~5.

P.~Flajolet, X.~Gourdon and P.~Dumas, ``Mellin transforms and asymptotics: harmonic sums'', \emph{Theoret.\ Comput.\ Sci.} 144 (1995), 3--58, is the reference for Section~\ref{sec:mellin}; P.~Flajolet
and G.~N.~Martin, ``Probabilistic counting algorithms for data base applications'', \emph{J.\ Comput.\ System Sci.} 31 (1985), 182--209; P.~Flajolet, \'E.~Fusy, O.~Gandouet and F.~Meunier,
``HyperLogLog: the analysis of a near-optimal cardinality estimation algorithm'', \emph{DMTCS Proceedings} AH (2007), 127--146. P.~Flajolet and R.~Sedgewick, \emph{Analytic Combinatorics}
(Cambridge, 2009), is the grand synthesis.
